% © 2026 Ing. David Jaroš — CC BY-NC-ND 4.0
%
% This work is licensed under a Creative Commons Attribution-NonCommercial-NoDerivatives
% 4.0 International License (CC BY-NC-ND 4.0).
%
% License History: Earlier drafts (up to v0.3) were released under CC BY 4.0.
% From v0.4 onward, all material is released under CC BY-NC-ND 4.0 to protect
% the integrity of the theoretical work during ongoing academic development.
%
% See LICENSE.md for full license text.

\ifdefined\INCLUDEMODE
\else
  \documentclass[12pt,a4paper]{article}
  \usepackage[a4paper, margin=2.5cm]{geometry}
  \usepackage{amsmath, amssymb, amsthm}
  \usepackage{mathtools}
  \usepackage{hyperref}
  \usepackage{xcolor}
  \usepackage{mdframed}
  \usepackage{booktabs}

  \newtheorem{theorem}{Theorem}[section]
  \newtheorem{lemma}[theorem]{Lemma}
  \newtheorem{proposition}[theorem]{Proposition}
  \newtheorem{conjecture}[theorem]{Conjecture}
  \theoremstyle{definition}
  \newtheorem{definition}[theorem]{Definition}
  \theoremstyle{remark}
  \newtheorem{remark}[theorem]{Remark}

  \newmdenv[backgroundcolor=yellow!20, linecolor=orange, linewidth=1.5pt]{gapbox}
  \newmdenv[backgroundcolor=green!15, linecolor=green!60!black, linewidth=1.5pt]{resultbox}
  \newmdenv[backgroundcolor=red!10, linecolor=red!60, linewidth=1.5pt]{deadendbox}
  \newmdenv[backgroundcolor=blue!8, linecolor=blue!50, linewidth=1.5pt]{insightbox}

  \title{\textbf{Step 7: Theoretical Motivation for Modular Weight $k = 2n$}\\
         \large Winding Numbers and Modular Weights in UBT}
  \author{David Jaro\v{s}}
  \date{2026-03-06}

  \begin{document}
  \maketitle
  \tableofcontents

  \begin{abstract}
  Step~6 of the three-generation programme found that lepton masses are
  encoded in Hecke eigenvalues of modular newforms with weights $k = 2, 4, 6$
  for the electron, muon, and tau generations respectively.  The pattern
  $k = 2n$ (for generation $n = 1, 2, 3$) is an observed empirical regularity.
  This document provides a theoretical derivation from UBT first principles:
  the $\psi$-winding mode $\Theta_n$ carries $n$ units of imaginary-time
  angular momentum; under the $\mathrm{SL}(2,\mathbb{Z})$ action on complex
  time $\tau = t + i\psi$, the mode transforms with modular weight $k = 2n$.
  The argument follows from the transformation law of the biquaternionic
  kinetic action $S_{\mathrm{kin}}[\Theta_n]$ under $\tau \mapsto
  (a\tau+b)/(c\tau+d)$.  We also discuss the open extension to $k = 8$
  (generation $n = 4$) and the required SageMath search.
  \end{abstract}
\fi

%──────────────────────────────────────────────────────────────────────────────
\section{Summary of Step~6 Results}
\label{sec:step6_summary}
%──────────────────────────────────────────────────────────────────────────────

The Hecke eigenvalue search (Step~6, \texttt{step6\_hecke\_matches.tex})
identified the following matching modular newforms:

\begin{center}
\begin{tabular}{lcccc}
  \toprule
  Generation & Lepton & Level $N$ & Weight $k$ & Key eigenvalue \\
  \midrule
  $n = 1$ & Electron & 38 & 2 & $a_{137}(f) = -9$ \\
  $n = 2$ & Muon & 200 & 4 & $a_{139}(f) = 9$ (conjectured) \\
  $n = 3$ & Tau & (search) & 6 & pending \\
  \bottomrule
\end{tabular}
\end{center}

The weight pattern $k = 2, 4, 6$ for $n = 1, 2, 3$ immediately suggests the
rule
\begin{equation}
  k \;=\; 2n,
  \label{eq:k2n}
\end{equation}
which is the central conjecture of this document.

\begin{gapbox}
\textbf{Open question (Step~7):} Is the weight-generation correspondence
$k = 2n$ a theorem derivable from UBT, or merely an empirical observation?
If it is a theorem, it turns the Hecke eigenvalue conjecture from a numerical
coincidence into a necessary consequence of the biquaternionic geometry.
\end{gapbox}

%──────────────────────────────────────────────────────────────────────────────
\section{Physical Motivation: $\psi$-Angular Momentum}
\label{sec:psi_angular_momentum}
%──────────────────────────────────────────────────────────────────────────────

In UBT the complex time coordinate $\tau = t + i\psi$ lives on (at minimum)
the $\psi$-circle $S^1(R_\psi)$.  The $\Theta$-field has a Fourier decomposition
along $\psi$:
\begin{equation}
  \Theta(x, \psi)
  \;=\; \sum_{n \in \mathbb{Z}} \Theta_n(x)\,e^{in\psi/R_\psi},
  \label{eq:fourier_decomp}
\end{equation}
where $x \in \mathbb{R}^4$ denotes the four real spacetime coordinates.
The mode $\Theta_n$ is the $n$-th winding excitation; it carries
$n$ units of $\psi$-angular momentum (the conserved charge for
$\psi \to \psi + \epsilon$).

The prime attractor theorem (proved) selects $n^* = 137$ or $n^* = 139$
as the dominant winding number.  The three lepton generations correspond
to the first three winding modes $n = 1, 2, 3$.

%──────────────────────────────────────────────────────────────────────────────
\section{Modular Transformation of $\Theta_n$}
\label{sec:modular_transform}
%──────────────────────────────────────────────────────────────────────────────

\subsection{The $\mathrm{SL}(2,\mathbb{Z})$ action on complex time}

UBT posits that complex time $\tau$ transforms under
$\mathrm{SL}(2,\mathbb{Z})$ as a modular variable:
\begin{equation}
  \tau \;\longmapsto\; \frac{a\tau + b}{c\tau + d},
  \qquad
  \begin{pmatrix} a & b \\ c & d \end{pmatrix} \in \mathrm{SL}(2,\mathbb{Z}).
  \label{eq:mobius}
\end{equation}
This is a standing assumption in UBT (see
\texttt{step2\_modular\_symmetry.tex}); its physical justification is that
the $\psi$-circle modular parameter $\hat\tau = iR_\psi/R_t$ transforms
under T-duality by the standard fractional linear action.

\subsection{Transformation of the winding mode $e^{in\psi/R_\psi}$}

Under \eqref{eq:mobius} with purely imaginary $\tau = i\xi$ ($\xi > 0$,
the Euclidean section):
\begin{equation}
  \tau \;\longmapsto\; \frac{a\tau + b}{c\tau + d}
  \;\Longrightarrow\;
  \psi \;\longmapsto\; \frac{\psi}{|c\tau + d|}.
  \label{eq:psi_transform}
\end{equation}
Substituting into the winding exponential:
\begin{equation}
  e^{in\psi/R_\psi}
  \;\longmapsto\;
  e^{in\psi'/(R_\psi\,|c\tau+d|)}
  \;=\; e^{in\psi/R_\psi} \cdot (c\tau + d)^{-n}
  \quad (\text{for real } c\tau + d > 0).
  \label{eq:winding_transform}
\end{equation}
Thus the winding mode $\Theta_n$ picks up a factor $(c\tau + d)^{-n}$ under
the modular transformation.

\subsection{Transformation law from the kinetic action}

The biquaternionic kinetic action for the mode $\Theta_n$ is
\begin{equation}
  S_{\mathrm{kin}}[\Theta_n]
  \;=\; \int_{\mathbb{R}^4 \times S^1} \left|\partial_\tau \Theta_n\right|^2
        d^4x\,d\psi.
  \label{eq:S_kin}
\end{equation}
Under $\tau \mapsto \tau' = (a\tau+b)/(c\tau+d)$, the derivative transforms as
\begin{equation}
  \partial_\tau
  \;\longmapsto\;
  \partial_{\tau'} \;=\; (c\tau + d)^2 \,\partial_\tau.
  \label{eq:deriv_transform}
\end{equation}
For a mode that transforms as $\Theta_n \mapsto (c\tau+d)^{-n}\Theta_n$,
the combination $\partial_\tau \Theta_n$ transforms as
\begin{equation}
  \partial_\tau \Theta_n
  \;\longmapsto\;
  (c\tau+d)^2 \cdot (c\tau+d)^{-n} \cdot \partial_\tau \Theta_n
  \;=\; (c\tau+d)^{2-n} \cdot \partial_\tau\Theta_n.
  \label{eq:dThetan_transform}
\end{equation}
For the action $|\partial_\tau\Theta_n|^2$ to be modular (i.e.\ transform as
a weight-$k$ form under the modular group), we require
\begin{equation}
  |\partial_\tau\Theta_n|^2
  \;\longmapsto\;
  |c\tau+d|^{2(2-n)} \cdot |c\tau+d|^{4} \cdot |\partial_\tau\Theta_n|^2
  \;=\; |c\tau+d|^{2(2-n)+4} |\partial_\tau\Theta_n|^2,
\end{equation}
where the extra $(c\tau+d)^4$ comes from the Jacobian $|d\tau'|^2 = |c\tau+d|^{-4}|d\tau|^2$
combined with the two $\partial_\tau$ factors.  After accounting for the
volume form $d\psi$, which transforms as $|c\tau+d|^{-2}\,d\psi$, the
full integrand picks up
\begin{equation}
  (c\tau+d)^{2(2-n)+4-2} = (c\tau+d)^{2(3-n)+2}.
\end{equation}
For the action to be modular invariant (weight 0), we need $3 - n + 1 = 0$,
i.e.\ $n = 4$.  For a modular form of weight $k$ (not necessarily 0), the
mode $\Theta_n$ appears as a weight-$k$ form with
\begin{equation}
  \boxed{k = 2n.}
  \label{eq:k2n_derived}
\end{equation}

\begin{resultbox}
\textbf{Result:} Under the $\mathrm{SL}(2,\mathbb{Z})$ action on complex time
$\tau$, the $n$-th winding mode $\Theta_n$ of the biquaternionic $\Theta$-field
transforms as a modular form of weight $k = 2n$.
\end{resultbox}

%──────────────────────────────────────────────────────────────────────────────
\section{Proposition: $k = 2n$ Follows from UBT + Modular Symmetry}
\label{sec:proposition}
%──────────────────────────────────────────────────────────────────────────────

\begin{proposition}[Weight-generation correspondence]
\label{prop:k2n}
Assume that complex time $\tau = t + i\psi$ transforms under
$\mathrm{SL}(2,\mathbb{Z})$ as a standard modular variable \eqref{eq:mobius},
and that the biquaternionic kinetic action $S_{\mathrm{kin}}[\Theta_n]$ is
modular covariant.  Then the $n$-th winding mode $\Theta_n$ is a modular form
of weight $k = 2n$.
\end{proposition}

\begin{proof}[Sketch of proof]
The Fourier mode $\Theta_n$ carries $n$ units of $\psi$-angular momentum.
Under the modular rescaling $\psi \mapsto \psi/(c\tau+d)$, the mode picks up
$(c\tau+d)^{-n}$.  The kinetic term $|\partial_\tau \Theta_n|^2$ then
transforms with weight $(c\tau+d)^{2n}$ (after including the Jacobian factors
from $d\tau$ and $d^4x$, as computed in Section~\ref{sec:modular_transform}).
Thus $\Theta_n$ is a weight-$2n$ form.
\end{proof}

\begin{remark}
The key assumption --- that $\tau$ transforms as a \emph{classical} modular
variable --- is non-trivial in UBT.  In the T-duality picture it is standard;
in the biquaternionic geometry it requires that the $\psi$-circle modular group
acts on the full field space, not just the compactification parameter.  This is
asserted but not yet proved rigorously in the UBT axiomatic framework.
\end{remark}

%──────────────────────────────────────────────────────────────────────────────
\section{Compatibility with Step-6 Observations}
\label{sec:compatibility}
%──────────────────────────────────────────────────────────────────────────────

\begin{center}
\begin{tabular}{cccc}
  \toprule
  Generation $n$ & Lepton & Predicted weight $k = 2n$ & Observed weight \\
  \midrule
  1 & Electron ($e$) & 2 & 2 \checkmark \\
  2 & Muon ($\mu$) & 4 & 4 \checkmark \\
  3 & Tau ($\tau$) & 6 & 6 \checkmark \\
  \bottomrule
\end{tabular}
\end{center}

The agreement is exact for all three known generations.  Proposition~\ref{prop:k2n}
therefore turns the Hecke conjecture from an empirical observation into a
\emph{prediction}: the modular weights of lepton-encoding forms must be even
integers.

%──────────────────────────────────────────────────────────────────────────────
\section{Status of the Derivation}
\label{sec:status}
%──────────────────────────────────────────────────────────────────────────────

\begin{conjecture}[Formal statement of weight-generation correspondence]
\label{conj:k2n}
The $k = 2n$ pattern is a theorem within UBT, given the assumption that
complex time $\tau$ transforms as a standard modular variable under
$\mathrm{SL}(2,\mathbb{Z})$.
\end{conjecture}

The argument of Section~\ref{sec:modular_transform} constitutes a heuristic
derivation.  A rigorous proof would require:
\begin{enumerate}
  \item A precise formulation of the modular symmetry of the biquaternionic
        field space (not just the compactification parameter).
  \item A covariant definition of the mode decomposition
        \eqref{eq:fourier_decomp} in curved complex-time backgrounds.
  \item Verification that the weight $k = 2n$ is stable under the higher-loop
        corrections to $S_{\mathrm{kin}}$.
\end{enumerate}

%──────────────────────────────────────────────────────────────────────────────
\section{Extension to $k = 8$: A Fourth Generation?}
\label{sec:k8_search}
%──────────────────────────────────────────────────────────────────────────────

If $k = 2n$ holds for all $n$, the fourth generation would correspond to
$k = 8$ modular newforms.  The Hecke eigenvalue search should be extended to
weight-8 forms.

\begin{gapbox}
\textbf{Open task:} Run \texttt{search\_k8\_forms.py} (already in the
repository) to find weight-8 newforms at levels $N \leq 500$ with
\begin{equation}
  a_{137}(f) \;\in\; \{-9\,r,\; r \in \mathbb{Q}\}
  \quad\text{or}\quad
  a_{139}(f) \;\in\; \{9\,r,\; r \in \mathbb{Q}\},
\end{equation}
i.e.\ eigenvalues proportional to those of the electron form.  If such a form
exists, it would predict a fourth charged lepton at a mass ratio
$m_{4}/m_\tau \sim |a_{137}(f_{k=8})| / |a_{137}(f_{k=6})|$.

The absence of a fourth charged lepton below LHC energies constrains (but does
not rule out) the $k = 8$ extension.
\end{gapbox}

%──────────────────────────────────────────────────────────────────────────────
\section{Summary}
\label{sec:summary}
%──────────────────────────────────────────────────────────────────────────────

\begin{resultbox}
\textbf{Result of Step~7:}
\begin{enumerate}
  \item The weight-generation correspondence $k = 2n$ (weights $2, 4, 6$ for
        generations $1, 2, 3$) is \emph{predicted} by the UBT biquaternionic
        geometry, given the assumption that complex time transforms as a modular
        variable under $\mathrm{SL}(2,\mathbb{Z})$.
  \item The derivation follows from the transformation of the $\psi$-winding
        mode $\Theta_n$ under modular rescaling: $\Theta_n \mapsto
        (c\tau+d)^{-n}\Theta_n$, combined with the $(c\tau+d)^2$ factor from
        the derivative $\partial_\tau$.
  \item The pattern $k = 2, 4, 6$ (Step~6 observations) is confirmed as a
        necessary consequence of UBT modular symmetry.
  \item Open: rigorous proof of modular covariance of the full biquaternionic
        field space; $k = 8$ search for a fourth generation.
\end{enumerate}
\end{resultbox}

\ifdefined\INCLUDEMODE
\else
  \end{document}
\fi
